% paper/sections/11a_force_balance.tex
%
% §11.1  力学的バランス検証（Force Balance Verification）
% ═══════════════════════════════════════════════════════════════════════════════
\subsection{力学的バランス検証}
\label{sec:force_balance}
% ═══════════════════════════════════════════════════════════════════════════════

NS ソルバの最も基本的な物理的整合性は，
\textbf{静止流体における力の釣合が数値的に維持されること}である．
本節では静水圧均衡と表面張力均衡の 2 種類のテストにより検証する．

% ─────────────────────────────────────────────────────────────────────────────
\subsubsection{静水圧均衡}
\label{sec:hydrostatic}
% ─────────────────────────────────────────────────────────────────────────────

\noindent\textbf{目的：}
Predictor--PPE--Corrector パイプラインが重力と圧力勾配の離散的釣合を
厳密に維持することを確認する．

\noindent\textbf{評価対象：}
NS ソルバ全体の離散力平衡（CCD 微分 + PPE 投影 + 速度補正）．

\noindent\textbf{手順とパラメタ：}
\begin{itemize}[noitemsep]
  \item 計算領域 $[0,1]^2$，単相流，壁面境界条件，重力 $\bm{g} = (0, -1)^\top$
  \item 初期条件 $\bu = \bm{0}$，$p = \rho g (1-y)$（静水圧から出発）
  \item 格子 $N = 32, 64, 128$，100 ステップ実行
  \item 合格基準 $\|\bu\|_\infty < 10^{-12}$
\end{itemize}

\noindent\textbf{結果：}

\begin{table}[htbp]
\centering
\caption{静水圧均衡テスト：100 ステップ後の速度・圧力誤差}
\label{tab:hydrostatic}
\renewcommand{\arraystretch}{1.3}
\begin{tabular}{rrrc}
\toprule
$N$ & $\|\bu\|_\infty$ & $|p - p_{\mathrm{hydro}}|_\infty$ & 判定 \\
\midrule
32  & $8.56 \times 10^{-16}$ & $4.89 \times 10^{-15}$ & 合格 \\
64  & $2.47 \times 10^{-15}$ & $8.11 \times 10^{-15}$ & 合格 \\
128 & $3.22 \times 10^{-15}$ & $9.99 \times 10^{-15}$ & 合格 \\
\bottomrule
\end{tabular}
\end{table}

\noindent\textbf{結論：}
速度・圧力ともに計算機精度水準（$\Ord{10^{-15}}$）に収まり，
重力と圧力勾配の離散的釣合が厳密に成立した．
CCD 微分演算子のフリーストリーム保存（§\ref{sec:gcl_verification}）の直接的帰結．

% ─────────────────────────────────────────────────────────────────────────────
\subsubsection{表面張力均衡と寄生流れ}
\label{sec:bf_static_droplet}
% ─────────────────────────────────────────────────────────────────────────────

\noindent\textbf{目的：}
静止液滴における Laplace 圧均衡（$\Delta p = \sigma/R$）を検証し，
CSF + 変密度 PPE の離散化が正しい圧力跳びを再現することを確認する．

\noindent\textbf{評価対象：}
CCD 曲率 + CSF 体積力 + 変密度 PPE（$\bnabla\cdot({\rho^{-1}}\bnabla p)$）の
単ステップ投影による Laplace 圧精度．

\noindent\textbf{手順とパラメタ：}
\begin{itemize}[noitemsep]
  \item 計算領域 $[0,1]^2$，壁面境界条件，重力なし
  \item 静止液滴 $R = 0.25$，$\rho_l/\rho_g = 1000$，$\mathit{We} = 1$
  \item CSF 体積力 $\bm{f}_\sigma = (1/\mathit{We})\,\kappa\,\bnabla\psi$（無次元形式では $\sigma$ は $We$ に吸収済み）
  \item 単ステップ投影：$\bu^* = \Delta t\,\bm{f}_\sigma / \rho$ → 変密度 PPE → Corrector
  \item PPE 右辺の発散は CCD 微分（$\Ord{h^6}$）で計算
  \item 格子 $N = 32, 64, 128, 256$，厳密解 $\Delta p = \sigma / R = 4.0$
\end{itemize}

\noindent\textbf{結果：}

\begin{table}[htbp]
\centering
\caption{静止液滴 Laplace 圧テスト：変密度 PPE + CSF（$\rho_l/\rho_g = 1000$）}
\label{tab:bf_laplace}
\renewcommand{\arraystretch}{1.3}
\begin{tabular}{rrrrr}
\toprule
$N$ & $\Delta p$ & 相対誤差 & $\|\bu_{\mathrm{para}}\|_\infty$ & $\Delta p$ 次数 \\
\midrule
 32 & $-7.80$ & $2.9 \times 10^{0}$  & $6.8 \times 10^{-1}$ & --- \\
 64 & $3.69$  & $7.9 \times 10^{-2}$ & $4.4 \times 10^{-2}$ & $5.23$ \\
128 & $3.98$  & $6.1 \times 10^{-3}$ & $1.5 \times 10^{-2}$ & $3.69$ \\
256 & $3.99$  & $2.2 \times 10^{-3}$ & $2.4 \times 10^{-1}$ & $1.45$ \\
\bottomrule
\end{tabular}
\end{table}

\noindent\textbf{結論：}
$N = 64, 128, 256$ の 3 水準で $\Delta p$ の相対誤差は単調減少し，
収束次数は $5.23 \to 3.69 \to 1.45$ と漸近的に CSF モデル誤差 $\Ord{h^2}$ に近づく%
\footnote{$N = 32$（$R/h = 8$）は界面幅 $\varepsilon$ が液滴半径 $R$ に対して過大となり，
$\delta_\varepsilon$ 正則化の仮定が破綻するため漸近領域外にある．}．
$N = 256$ で $\Delta p = 3.99$（相対誤差 $0.22\%$）であり，
CCD-PPE の高精度離散化が CSF $\Ord{h^2}$ の枠内で十分に機能している．

寄生流れ振幅は $N = 64 \to 128$ で減少するが，$N = 256$ で増大する．
これは PPE ソルバ（FD $\Ord{h^2}$）と Corrector の CCD 圧力勾配との間の
Balanced-Force 不整合が高解像度で顕在化するためである（§~\ref{sec:balanced_force}）．
この不整合が多ステップ時間積分において累積するかを
§~\ref{sec:advection_pressure_coupling} の二相流時間積分テストで確認する．
寄生流れ振幅の格子収束（Balanced-Force 条件下での $\Ord{h^2}$）は，
時間積分による定常化が必要であり，
第~\ref{sec:validation}章の静止液滴ベンチマーク（§\ref{sec:val_static_drop}）にて検証する．

\paragraph{CCD $p''$ 係数相殺補正の効果}

§~\ref{sec:rc_high_order} の CCD $p''$ による RC ブラケット係数相殺補正を
多ステップ静止液滴（400 ステップ，$\rho_l/\rho_g = 2$，$\mathit{We} = 10$，
Rhie--Chow Balanced-Force 拡張有効）に適用した結果を示す．

\begin{table}[htbp]
\centering
\caption{CCD $p''$ 係数相殺補正の寄生流れへの効果（400 ステップ後）}
\label{tab:rc_d2fd_droplet}
\renewcommand{\arraystretch}{1.3}
\begin{tabular}{clccr}
\toprule
$N$ & RC 方式 & $\|\bu_{\mathrm{para}}\|_\infty$ & $\Delta p$ 相対誤差 & 寄生流れ低減率 \\
\midrule
32 & 標準 $\Ord{h^2}$       & $2.51 \times 10^{-2}$ & $0.58\%$ & --- \\
32 & $p''$ 補正 $\Ord{h^4}$  & $2.49 \times 10^{-2}$ & $0.42\%$ & $\mathbf{-0.8\%}$ \\
\midrule
64 & 標準 $\Ord{h^2}$       & $5.49 \times 10^{-4}$ & $1.23\%$ & --- \\
64 & $p''$ 補正 $\Ord{h^4}$  & $4.24 \times 10^{-4}$ & $1.23\%$ & $\mathbf{-23\%}$ \\
\bottomrule
\end{tabular}
\end{table}

$N = 64$ では $p''$ 係数相殺補正により寄生流れ振幅が 23\% 低減し，
$\Delta p$ 精度は同等に維持される．
$N = 32$ でも $\Delta p$ が改善（$0.58\% \to 0.42\%$）しており，
標準 RC の構造を保持したまま補正を加算する方式が壁境界でも安定に動作することを示している．

CCD が $p'$ と同時に返す $p''$ のみを使用するため，追加の CCD 呼び出しは不要であり，
従来の FD ベースの RC では原理的に得られない\textbf{高次微分情報の活用}が，
追加コストゼロの精度改善を可能にしている．
